\KashidaOff
% در این فایل، عنوان پایان‌نامه، مشخصات خود، متن تقدیمی‌، ستایش، سپاس‌گزاری و چکیده پایان‌نامه را به فارسی، وارد کنید.
% توجه داشته باشید که جدول حاوی مشخصات پایان‌نامه/رساله و همچنین، مشخصات داخل آن، به طور خودکار، درج می‌شود.
%%%%%%%%%%%%%%%%%%%%%%%%%%%%%%%%%%%%
% دانشگاه خود را وارد کنید
\university{دانشگاه فردوسی}
% دانشکده، آموزشکده و یا پژوهشکده  خود را وارد کنید
\faculty{دانشکده مهندسی}
% گروه آموزشی خود را وارد کنید
\department{مهندسی کامپیوتر}
% گروه آموزشی خود را وارد کنید
\subject{نرم افزار}

% گرایش خود را وارد کنید
%\field{گرایش‌تان}
% عنوان پایان‌نامه را وارد کنید
\title{مدل محاسباتی همگام فدرالی براي شبکه هاي تطبیقی توزیع شده}
% نام استاد(ان) راهنما را وارد کنید
\firstsupervisor{دکتر عبدالرضا سوادی}
%\secondsupervisor{استاد راهنمای دوم}
% نام استاد(دان) مشاور را وارد کنید. چنانچه استاد مشاور ندارید، دستور پایین را غیرفعال کنید.
\firstadvisor{دکتر هادی صدوقی یزدی}
%\secondadvisor{استاد مشاور دوم}
% نام پژوهشگر را وارد کنید
\name{فاطمه}
% نام خانوادگی پژوهشگر را وارد کنید
\surname{بارانی}
% تاریخ پایان‌نامه را وارد کنید
\thesisdate{1403}
\okdate{1400/00/00}
\defencedate{1400/11/11}
% کلمات کلیدی پایان‌نامه را وارد کنید
\keywords{مدل محاسباتی موازی، یادگیری توزیع شده، شبکه های تطبیقی توزیع شده}
% چکیده پایان‌نامه را وارد کنید
\fa-abstract{
حداکثر در حجمی معادل با 250 تا 300 کلمه تهيه شده و شامل بيان مختصر مسئله مورد بررسی، روش تحقيق و مراحل بکار گرفته شده برای کسب و جمع آوری اطلاعات، نحوه تجزيه و تحليل و نتيجه کلی می باشد. خواننده با مطالعه چکيده بايد تشخيص دهد که رساله موجود دربرگيرنده مطالب مورد علاقه وی می باشد يا خير؟ تاريخچه و سابقه موضوع در اين قسمت ذکر نشده، بلکه در مقدمه رساله توضيح داده می شود. چکيده در يک صفحه مجزا قبل از فهرست مطالب قرار می گيرد . در بالای آن به فاصله دو سطر از حاشيه بالای صفحه در ميانه سطر عنوان پايان نامه نوشته می شود. در انتهای چکيده ميتواند کلمات کليدی مورد استفاده در پايان نامه به تعداد 5-4 واژه اضافه شود. \\
دوستان شما در این رساله سعی دارند تا شما را با یک قالب پایان‌نامه/ رساله آشنا کنند. شما با توجه به همین بسته موجود (FThesis) و راهنمایی‌های ارائه شده در این نمونه رساله خواهید توانست با اندکی دقت ضمن یاد گرفتن اصول فنی نوشتن تحت\TeX با نگارش فنی نیز آشنا شوید، لازم به ذکر است که قالب حاضر به طور اختصاصی استانداردهای دانشگاه فردوسی را پشتیبانی می‌کند. \\
لازم به ذکر است که ما تعداد کلمات در چکیده به طور رسمی دارای محدودیت است از این‌رو ما نیز با توجه به آن فضای مربوط به چکیده‌مان را تنظیم کردیم.
}
\ftitle
 % پایان‌نامه خود را تقدیم کنید!
\begin{acknowledgementpage}

\vspace{4cm}

{\nastaliq
{\LARGE
 تقدیم به همه آنهایی که 
\vspace{1.5cm}

\hspace{3cm}
می خواهند بیشتر بدانند
}}
\end{acknowledgementpage}
\newpage
%%%%%%%%%%%%%%%%%%%%%%%%%%%%%%%%%%%%
% ستایش
\baselineskip=.750cm
\thispagestyle{empty}
% \KashidaOff
{\nastaliq
خدایا...%
\RTLfootnote{مناجاتی از دکتر علی شریعتی.}
}\\
\vspace{.5cm}\\

به من زیستنی عطا کن که در لحظه مرگ، بر بی‌ثمری لحظه‌ای که برای زیستن گذشته است، حسرت نخورم   و مُردنی عطا کن که بر بیهودگیش، سوگوار نباشم. بگذار تا آن را، خود انتخاب کنم، اما آنچنان که تو دوست می‌داری.

تو می‌دانی و همه می‌دانند که شکنجه دیدن بخاطر تو، زندانی کشیدن بخاطر تو و رنج بردن به پای تو تنها لذت بزرگ زندگی من است، از شادی توست که من در دل می‌خندم، از امید رهایی توست که برق امید در چشمان خسته‌ام می‌درخشد و از خوشبختی توست که هوای پاک سعادت را در ریه‌هایم احساس می‌کنم. نمی‌توانم خوب حرف بزنم. نیروی شگفتی را که در زیر کلمات ساده و جمله‌های ضعیف و افتاده، پنهان کرده‌ام دریاب، دریاب.

تو می‌دانی و همه می‌دانند که زندگی از تحمیل لبخندی بر لبان من، از آوردن برق امیدی در نگاه من، از برانگیختن موج شعفی در دل من، عاجز است.

تو، چگونه زیستن را به من بیاموز، چگونه مردن را خود خواهم آموخت.

به من توفیق تلاش در شکست، صبر در نومیدی، رفتن بی‌همراه، جهاد بی‌سلاح، کار بی‌پاداش، فداکاری در سکوت، دین بی‌دنیا، مذهب بی‌عوام، عظمت بی‌نام، خدمت بی‌نان، ایمان بی‌ریا، خوبی بی‌نمود، گستاخی بی‌خامی، قناعت بی‌غرور، عشق بی‌هوس، تنهایی در انبوه جمعیت، و دوست داشتن بی‌آنکه دوست بداند،  روزی کن. 

\vspace{1.5cm}
{\nastaliq
\hspace{1cm}
اگر تنها‌ترین تنها شوم، باز خدا هست
\vspace{.8cm}

\hspace{4.7cm}
او جانشین همه نداشتن‌هاست...
}
%%%%%%%%%%%%%%%%%%%%%%%%%%%%%%%%%%%%
\newpage\thispagestyle{empty}
% سپاس‌گزاری
{\nastaliq
سپاس‌گزاری...
}
\\[2cm]
سپاس خداوندگار حکیم را که با لطف بی‌کران خود، آدمی را زیور عقل آراست.


در آغاز وظیفه‌  خود  می‌دانم از زحمات بی‌دریغ استاد  راهنمای خود،  جناب آقای دکتر  سوادی، صمیمانه تشکر
و  قدردانی کنم  که قطعاً بدون 
راهنمایی‌های ارزنده‌  ایشان، این مجموعه  به انجام  نمی‌رسید.

از جناب  آقای  دکتر صدوقی یزدی که زحمت  مطالعه و مشاوره‌  این رساله
را تقبل  فرمودند و
در آماده سازی  این رساله، به نحو احسن اینجانب را مورد راهنمایی قرار دادند، کمال امتنان را دارم.

 در پایان، بوسه می‌زنم بر دستان خداوندگاران مهر و مهربانی، پدر و مادر عزیزم و بعد از خدا، ستایش می‌کنم وجود مقدس‌شان را و تشکر می‌کنم از همسر و فرزندانم که با صبوری مرا در این مسیر حمایت کردند.
% با استفاده از دستور زیر، امضای شما، به طور خودکار، درج می‌شود
\signature 
\newpage\clearpage
\KashidaOn